% !TeX root = ../main.tex
\newcommand{\SplitSeed}{20261103}
\newcommand{\LabeledDuplicateExcess}{609}
\newcommand{\FeatureDuplicateGroups}{354}
\newcommand{\FeatureDuplicateRows}{1,000}
\newcommand{\CrossSplitFeatureGroups}{246}


以 \texttt{SeriousDlqin2yrs} 为分层标签，固定随机种子为 \SplitSeed，
将全部 150,000 条记录按 $60\%/20\%/20\%$ 划分。
先分出 $60\%$ 的训练集，再将剩余 $40\%$ 等分为验证集和测试集；
两次均使用同一随机种子并按标签分层。
各子集规模与坏样本率见表~\ref{tab:sample-split}。
分层使各部分的事件比例接近总体的 $6.684\%$。

\begin{table}[htbp]
  \centering
  \caption{固定分层随机划分的样本构成（坏样本为正类）}
  \label{tab:sample-split}
  \small
  \begin{tabular}{lrrrrr}
\toprule
子集 & 样本数 & 好样本数 & 坏样本数 & 坏样本率 & 样本占比 \\
\midrule
训练集 & 90,000 & 83,984 & 6,016 & 6.684\% & 60\% \\
验证集 & 30,000 & 27,995 & 2,005 & 6.683\% & 20\% \\
测试集 & 30,000 & 27,995 & 2,005 & 6.683\% & 20\% \\
\midrule
合计 & 150,000 & 139,974 & 10,026 & 6.684\% & 100\% \\
\bottomrule
\end{tabular}

\end{table}

\paragraph{重复与同源记录核查}
原始行号唯一，但仅作为追踪键。
排除行号后，按十个特征和标签检查完全重复记录，发现
\LabeledDuplicateExcess{} 条相对于各组首条的重复行；
仅按十个特征（含缺失位置）匹配，有 \FeatureDuplicateGroups{} 组、
共 \FeatureDuplicateRows{} 条记录，其中 \CrossSplitFeatureGroups{} 组分布在多个子集。
因缺少客户主键（名称或者用户id等唯一编码信息）和逐户日期，无法确认这些记录是否来自同一借款人，
故本次保留原始全部行，按指定比例进行记录层面的分层留出，即以每行为划分单位，
在好、坏样本中分别按比例随机划分训练、验证和测试集，不要求同特征记录进入同一集合。


\paragraph{各部分用途}
训练集用于估计清洗参数、分箱切点、档位合并规则、WOE、IV 以及模型系数；
验证集用于比较和选择模型方案、概率校准方案与决策阈值；
待数据处理规则、模型及阈值全部冻结后，测试集仅承担最终评价，不据其结果回调方案。
若在训练集内部实施交叉验证，则每一折的分箱、WOE 编码、特征筛选及其他数据驱动处理
均须仅在该折训练部分重新学习，折内验证部分只应用已学得的规则，避免信息泄漏。
